\chapter{Spines and products of stationary cones}\label{sec:cone-spines-products}
Let $C$ be a nonzero stationary integral $m$-cone on $\mathbb R^N$, where $1\leq m<N$, and write $\mu=\|C\|$ and $\Theta(y)=\Theta^m(\mu,y)$. We develop the argument of \cite{Minter}, Lemma 1.5, pp.~13--15. A cone is invariant under dilations as a varifold, including the tangential Jacobian and the recorded planes.

\begin{proposition}[Density under scaling and at infinity]\label{prop:cone-density-infinity}
For every $y\in\mathbb R^N$ and $s>0$,
\[
 \Theta(sy)=\Theta(y),\qquad
 \lim_{R\to\infty}\frac{\mu(B_R(y))}{\omega_mR^m}=\Theta(0).
\]
Consequently $\Theta(y)\leq\Theta(0)$ for every $y$.
\end{proposition}
\begin{proof}
Homogeneity gives $\mu(B_r(sy))=s^m\mu(B_{r/s}(y))$. Divide by $\omega_mr^m$ and let $r\downarrow0$. For $R>|y|$, the inclusions
\[
 B_{R-|y|}(0)\subset B_R(y)\subset B_{R+|y|}(0)
\]
and $\mu(B_t(0))=\omega_mt^m\Theta(0)$ squeeze the ratio at $y$ between $\Theta(0)(1-|y|/R)^m$ and $\Theta(0)(1+|y|/R)^m$. Let $R\to\infty$. Stationary monotonicity, Theorem \ref{thm:stationary-monotonicity}, bounds the density at $y$ by this limit.
\end{proof}

\begin{proposition}[Maximal density gives another cone vertex]\label{prop:maximal-density-vertex}
For $y\in\mathbb R^N$, $\Theta(y)=\Theta(0)$ if and only if $C$ is invariant under every homothety $h_{y,s}(x)=y+s(x-y)$, $s>0$.
\end{proposition}
\begin{proof}
If the densities agree, the nondecreasing mass ratio centered at $y$ has equal limits at zero and infinity by Proposition \ref{prop:cone-density-infinity}. It is therefore constant. Translate the origin to $y$ and apply Theorem \ref{thm:constant-ratio-cone}; using reciprocal scale gives precisely the stated homotheties. Conversely these homotheties make every centered ratio at $y$ constant, so its zero and infinite scale limits coincide, again by Proposition \ref{prop:cone-density-infinity}.
\end{proof}

\begin{lemma}[Two vertices generate a line of translations]\label{lem:two-cone-vertices}
If a varifold is invariant under $h_{0,s}$ and $h_{y,s}$ for every $s>0$, then it is invariant under $\tau_{ty}(x)=x+ty$ for every real $t$.
\end{lemma}
\begin{proof}
The composition $h_{y,s}\circ h_{0,1/s}$ equals $x\mapsto x+(1-s)y$. Varifold pushforwards compose: the scalar tangential Jacobians multiply, and here $s^m s^{-m}=1$. Thus every translation with $t<1$ preserves the varifold. The inverse of a preserving translation also preserves it, supplying all remaining values of $t$. This is an identity of full varifolds, not just their supports.
\end{proof}

\begin{proposition}[The spine is the linear set of maximal density]\label{prop:spine-maximal-density}
The spine from Definition \ref{def:cone-spine} satisfies
\[
 S(C)=\{y:\Theta(y)=\Theta(0)\}.
\]
It is a closed linear subspace of dimension at most $m$.
\end{proposition}
\begin{proof}
A preserving translation gives $\mu(B_r(y))=\mu(B_r(0))$, hence equality of densities. The reverse implication follows from Proposition \ref{prop:maximal-density-vertex} and Lemma \ref{lem:two-cone-vertices}. Preserving translations are closed under addition and inverses; the lemma gives all real multiples of any one of them. They form a linear subspace. Closedness also follows from continuity of translations against compactly supported continuous tests, with all supports in a common compact set when the translation vectors converge.

Write $d=\dim S(C)$. The mass $b=\mu(B_1(0))$ is positive, by conical homogeneity and nonzeroness. A lattice in $S(C)$ with separation three has at least $c_dR^d$ points in $B_R(0)$ for large $R$ if $d>0$. The unit balls at these points are disjoint, each has mass $b$, and all lie in $B_{R+1}(0)$. Thus $c_dbR^d\leq b(R+1)^m$, forcing $d\leq m$. The case $d=0$ needs no estimate.
\end{proof}

\begin{lemma}[Factoring a translation-invariant Radon measure]\label{lem:translation-measure-product}
Let $L$ be a $d$-dimensional linear subspace, $d\geq1$. If a Radon measure $\mu$ on $L^\perp\times L$ is invariant under translations in $L$, there is a unique Radon measure $\nu$ on $L^\perp$ such that
\[
 \mu=\nu\otimes\mathcal L^d.
\]
Here orthonormal coordinates fix the Euclidean volume measure on $L$.
\end{lemma}
\begin{proof}
For a bounded Borel set $A\subset L^\perp$, the measure $B\mapsto\mu(A\times B)$ is locally finite and translation invariant on $L$. Such a measure is a constant multiple of $\mathcal L^d$: partition half-open cubes into equal subcubes, use translation invariance to give them equal mass, and then approximate arbitrary boxes by rational grids. A bounded piece of a coordinate hyperplane has zero mass, since arbitrarily many disjoint translates fit in a fixed bounded box; thus grid boundaries cause no ambiguity. Uniqueness of locally finite measures from boxes extends the equality to all Borel sets.

Let $Q$ be a half-open unit cube in $L$ and define $\nu(A)=\mu(A\times Q)$. This is countably additive and finite on bounded sets, and hence is a Radon measure in Euclidean space. The preceding argument gives $\mu(A\times B)=\nu(A)\mathcal L^d(B)$ for bounded rectangles. Exhaustion and uniqueness of product measures prove the claim for all Borel sets, as well as uniqueness of $\nu$.
\end{proof}

\begin{lemma}[Spine directions lie in the recorded planes]\label{lem:spine-tangent-directions}
If $L\subset S(C)$, then $L\subset S$ for $C$-almost every $(x,S)$.
\end{lemma}
\begin{proof}
The constant mass ratio at zero and Corollary \ref{cor:stationary-radial-equality} give $x\in S$ almost everywhere. Each vector $e$ in a fixed basis of $L$ is another cone vertex by Proposition \ref{prop:spine-maximal-density}, so the same argument gives $x-e\in S$ almost everywhere. Intersect the finitely many full-measure sets and subtract to obtain $e\in S$ for every basis vector.
\end{proof}

\begin{proposition}[Integral product splitting along a proper spine subspace]\label{prop:cone-product-splitting}
Suppose $L\subset S(C)$ and $1\leq d=\dim L<m$. There is a nonzero stationary integral $(m-d)$-cone $C_0$ on $L^\perp$ such that
\[
 C=C_0\times |L|.
\]
The product records the plane $T_yC_0\oplus L$ and the multiplicity of $C_0$ at $(y,z)$. The rectifiable coarea input is \cite{Sim18}, Chapter 3, (2.8)--(2.9) and Remark 2.10, pp.~93--94; compare the conclusion of \cite{Minter}, Lemma 1.5.
\end{proposition}
\begin{proof}
Write $C=V(M,\theta)$. Lemmas \ref{lem:translation-measure-product} and \ref{lem:spine-tangent-directions} give $\mu=\nu\otimes\mathcal L^d$ and $L\subset T_xM$ almost everywhere. Thus the tangential Jacobian of the projection $p:M\to L$ is one. Weighted rectifiable coarea, Corollary \ref{cor:weighted-rectifiable-coarea}, expresses $\mu$ as the integral over $z\in L$ of the measures
\[
 \nu_z=\theta(\,\cdot\,,z)\mathcal H^{m-d}\llcorner\{y:(y,z)\in M\}.
\]
For almost every $z$, this is a locally finite integral rectifiable measure. Local finiteness and integrality follow on a countable compact exhaustion, with $\theta\geq1$; exceptional sets of zero weighted measure meet almost every fiber in zero weighted measure. Coarea and the product identity give $\int_B\nu_z(A)\,dz=\nu(A)\mathcal L^d(B)$. Apply this to a countable collection of bounded boxes generating the Borel sets of $L^\perp$. It follows that $\nu_z=\nu$ as measures for almost every $z$.

Choose one such fiber and let $C_0$ be its integral rectifiable lift in $L^\perp$. The area formula on product charts gives a rectifiable product $C_0\times|L|$ with weight $\nu\otimes\mathcal L^d$ and planes $T_yC_0\oplus L$. Lemma \ref{lem:rectifiable-lift-unique} therefore identifies this product with $C$; no invariance of individual exceptional fibers was assumed.

For $Y\in C_c^1(L^\perp;L^\perp)$ choose $\chi\in C_c^1(L)$ with $\int\chi=1$ and test with $X(y,z)=\chi(z)Y(y)$. Divergence in the $L$ directions is zero, and product integration gives $0=\delta C(X)=\delta C_0(Y)$. Finally homogeneity of $\mu$ gives $\nu(sA)=s^{m-d}\nu(A)$ by testing rectangles and scaling their $L$ factor. The dilated lift of $C_0$ has the same weight as $C_0$, so Lemma \ref{lem:rectifiable-lift-unique} proves conical invariance. Nonzeroness follows from the product identity.
\end{proof}

\begin{corollary}[A full-dimensional spine forces a multiple plane]\label{cor:full-spine-plane}
If $\dim S(C)=m$, then $C=q|S(C)|$ for a positive integer $q$.
\end{corollary}
\begin{proof}
Set $L=S(C)$. Lemma \ref{lem:translation-measure-product} gives $\mu=\nu\otimes\mathcal L^m$. Homogeneity now implies $\nu(sA)=\nu(A)$ for every $s>0$. The dyadic annuli $\{2^{-j-1}\leq|y|<2^{-j}\}$ all have equal $\nu$ mass; their disjoint union lies in a compact ball. Local finiteness forces that mass to be zero. Scaling and exhaustion show $\nu$ is supported at zero, so $\mu=q\mathcal H^m\llcorner L$ for a finite $q>0$. Rectifiable lift uniqueness identifies $C=q|L|$, and integrality makes $q$ an integer. Equivalently the transverse factor is the zero-dimensional varifold $q\delta_0$.
\end{proof}
